%% ma: L12-B14-0052
%% lop: 12
%% chuong: 5
%% bai: 14
%% dang: 12.14.1
%% loai: TN4
%% muc: NB
%% dap_an: "D"
%% nguon: "To_n_12_PT_M_t_Ph_ng.pdf (Đề kiểm tra đánh giá kiến thức, Chuyên đề Phương trình mặt phẳng, Nguyễn Anh Khôi - Lê Minh Khôi), Câu 6"
%% kien_thuc: "Bài 14 (phương trình tổng quát của mặt phẳng, vectơ pháp tuyến)"
%% trang_thai: chinh-thuc
%% ghi_chu: "Tài liệu gốc không kèm đáp án; đáp án do Claude giải và kiểm bằng sympy."
\begin{ex}
Trong không gian $Oxyz$, vectơ nào sau đây \textbf{không phải} là vectơ pháp tuyến của mặt phẳng $(P)\colon x+3y-5z+2=0$?
\choice{$\vec n=(-3;-9;15)$}{$\vec n=(-1;-3;5)$}{$\vec n=(2;6;-10)$}{\True $\vec n=(-2;-6;-10)$}
\loigiai{$(P)$ có vectơ pháp tuyến $(1;3;-5)$. Các vectơ A, B, C lần lượt bằng $-3$, $-1$, $2$ lần vectơ này; $(-2;-6;-10)$ không cùng phương. Chọn D.}
\end{ex}
